\section{Keterbatasan}
\label{sec:limits}
Setiap konfigurasi dijalankan satu kali; keluaran LLM bervariasi antar run~\cite{arya2026wizard}, dan perbedaan yang kami laporkan antar konfigurasi (sekitar 0.01 pada cosine) berada dalam variabilitas tersebut. Semua metrik kualitas bersifat otomatis; cosine lintas bahasa dan cosine terjemahan balik adalah proksi yang diketahui gagal menangkap kesalahan terminologi dan kefasihan~\cite{freitag2021mqm,moon2020roundtrip}, yang tepat merupakan area di mana penyunting bertindak---evaluasi manusia oleh pakar domain dwibahasa diperlukan untuk mengkuantifikasi manfaat tersebut. Satu pasangan bahasa, satu tipe dokumen, dan satu glosarium dipelajari. Perbandingan tersebut mengaburkan pilihan model dengan peran penyuntingan pasca: validator aktif di semua konfigurasi, sehingga hasil struktural mengisolasi mereka, tetapi satu putaran penyuntingan pasca oleh model \emph{sama} dengan draf tidak diuji. Model QA juga berfungsi sebagai model cadangan (fallback), sehingga kelemahan sistematis Gemma 3 27B akan memengaruhi dua peran. Eksperimen A berkaitan dengan satu SDK vendor dan satu tingkat akun pada satu titik waktu; perilaku dapat berubah, dan penyedia lain mungkin menolak model yang tidak dikenal. Tinjauan dua-kritikus dan pengonversi DOCX adalah bagian dari agen yang dirilis tetapi tidak dievaluasi di sini.

\section{Kesimpulan}
\label{sec:conclusion}
Kami mendokumentasikan bahwa sebuah vendor SDK dapat senyap melayani model yang berbeda dari yang diminta dan mengubah event penggunaan yang mengungkapkannya menjadi mekanisme provenans bagi agen riset. Di sekitarnya kami membangun paper-agent, yang merutekan tugas penanganan paper ke model berdasarkan peran dan memvalidasi setiap keluaran model secara deterministik. Dalam perbandingan terkontrol pada 146 unit LaTeX, sebuah pipeline tiga-model tidak mengungguli satu model tunggal pada kualitas rata-rata atau pada integritas struktur---keduanya berasal dari validator---tetapi pipeline itu memperbaiki kesalahan semantik yang tidak terlihat oleh metrik; para validator menangkap tiga kelas kesalahan yang tidak terdeteksi oleh model mana pun, termasuk penilai. Penelitian lanjutan adalah evaluasi manusia terhadap koreksi penyunting, sebuah filter terpelajar untuk review selektif, dan memperluas provenans model yang melayani ke SDK lain. Sebagai pemeriksaan akhir, agen menerjemahkan paper ini ke Bahasa Indonesia dengan pipeline yang sama (131 unit, 4{,}200 kata; 131/131 unit secara struktural utuh; kepatuhan glosarium 95.7\%; cosine 0.88 langsung dan 0.96 terjemahan balik; model yang melayani GPT-5 mini, Claude Haiku 4.5 dan Gemma 3 27B); kedua versi dikompilasi dari bibliografi yang sama. Kode, metrik per-unit dan kedua versi bahasa dirilis di \url{https://github.com/devnolife/paper-agent}.
